\renewcommand{\bibname}{Quellen und Anschlüsse im Projekt}
\begin{thebibliography}{99}
\addcontentsline{toc}{chapter}{Quellen und Anschlüsse im Projekt}
\small
\interlinepenalty=10000

\bibitem{project40}
Martin Kunze: \emph{Die Grundlagen der Mathematik}, Band~40,
\emph{Gruppen}.
Projektquelle: \path{Bd. 40 - Gruppen.tex}.
Definition~40.2.1.1 (Gruppe); Sätze~40.2.3.1, 40.2.4.1--2 und
40.2.5.1 (Inverse und Kürzung); Abschnitte~40.3.1--2 und
Satz~40.4.1.1 (Homomorphismen und Isomorphismen).
Die Untergruppen-, Quotienten-, Wirkungs- und Präsentationstheorie
dieses Leitfadens ergänzt den dort bisher ausgearbeiteten Bestand.

\bibitem{projectreconstruction}
Martin Kunze: \emph{Die Grundlagen der Mathematik}, Band~40,
Abschnitt~40.3.3; Ergänzungen \emph{Gruppenrekonstruktion --
Lesefassung} und \emph{Gruppenrekonstruktion -- Beweistabellen}.
Projektquellen: \path{tex/b40/reconstruction/}.
Sätze~40.3.3.2--4 und 40.3.3.10 zur Einermengenschicht,
Gruppenrekonstruktion und Absorption durch den ganzen Träger.
Die allgemeine Einheitengruppentheorie von Monoiden und die
stärkere Starrheitsfassung mit beliebiger Halbgruppe als
Vergleichsstruktur werden hier nicht zusätzlich entwickelt.

\bibitem{milne}
J. S. Milne: \emph{Group Theory}, Version~4.01, November 2025.
Kapitel~1: Grundbegriffe, Untergruppen, Homomorphismen, Nebenklassen,
Normalteiler und direkte Produkte; insbesondere
Sätze~1.45--1.47 und Korollar~1.48 zu Isomorphiesätzen und Korrespondenz.
Kapitel~2, S.~32--37: freie Gruppen und Präsentationen.
Kapitel~3: Automorphismen, charakteristische Untergruppen,
semidirekte Produkte und Erweiterungen.
Kapitel~4, S.~57--61: Wirkungen, Bahnen, Stabilisatoren und Nebenklassenaktionen.
Verwendet werden die allgemeinen Gruppenaussagen dieser Abschnitte.
Die universellen Produktformulierungen, das Faserprodukt und die
Kardinalfassungen der Indexformeln sind als elementare Folgerungen
dieser Grundlagen formuliert.
\url{https://www.jmilne.org/math/CourseNotes/GT.pdf}.

\bibitem{une}
University of New England: \emph{Abstract Algebra, PMTH332}.
Vorlesungsskript, 2014, Kapitel~1--2, insbesondere
Proposition~2.3, S.~29, zur Charakterisierung assoziativer
Gruppenverknüpfungen durch beidseitige Lösbarkeit.
\url{https://mcs.une.edu.au/~pmth332/Lectures/Lecture%20Notes%20(old).pdf}.

\bibitem{conradquot}
Keith Conrad: \emph{Quotient Groups}.
University of Connecticut, Vorlesungsnotizen.
Abschnitt~2, insbesondere Sätze~2.11--2.12:
Normalitätskriterien; Abschnitt~3: Quotientengruppe
und Repräsentantenunabhängigkeit.
\url{https://kconrad.math.uconn.edu/blurbs/grouptheory/quotientgroups.pdf}.

\bibitem{conradhom}
Keith Conrad: \emph{Homomorphisms}.
University of Connecticut, Vorlesungsnotizen.
Abschnitte~6--7, insbesondere Sätze~6.9, 7.10 und~7.11:
Untergruppenbilder, Untergruppenurbilder, Kerne und Fasern.
\url{https://kconrad.math.uconn.edu/blurbs/grouptheory/homomorphisms.pdf}.

\bibitem{conradconj}
Keith Conrad: \emph{Conjugation in a Group}.
University of Connecticut, Vorlesungsnotizen.
Satz~3.5 zur Konjugation als Automorphismus; Abschnitt~6
zu Konjugationsklassen und Zentralisatoren.
\url{https://kconrad.math.uconn.edu/blurbs/grouptheory/conjclass.pdf}.

\bibitem{conradsemi}
Keith Conrad: \emph{Semidirect Products}.
University of Connecticut, Vorlesungsnotizen.
Insbesondere Satz~4.1 zum inneren semidirekten Produkt.
\url{https://kconrad.math.uconn.edu/blurbs/grouptheory/semidirect-product.pdf}.

\bibitem{conradsplit}
Keith Conrad: \emph{Splitting of Short Exact Sequences for Groups}.
University of Connecticut, Vorlesungsnotizen.
Satz~3.3: homomorphe Schnitte und semidirekte Zerlegung.
\url{https://kconrad.math.uconn.edu/blurbs/grouptheory/splittinggp.pdf}.

\bibitem{mcmullen}
Curtis T. McMullen: \emph{Topology. Course Notes}.
Harvard University, Math~131, Herbst 2013.
Abschnitt~14, S.~75--79: freie Produkte beliebiger Gruppen,
universelle Eigenschaft und Normalform (Satz~14.6).
\url{https://legacy-www.math.harvard.edu/~ctm/papers/home/text/class/harvard/131/course/course.pdf}.

\bibitem{liutringali}
Shuolin Liu und Salvatore Tringali:
\emph{Power Semigroups and Two Rigidity Theorems for Groups}.
Preprint, arXiv:2606.01917v1, 1.~Juni 2026.
Einleitung und Abschnitt~2, insbesondere Bemerkung~2.2(2),
zur Erkennung von Einermengen; Satz~1.3 enthält eine stärkere
Starrheitsfassung. Der vorliegende Leitfaden verwendet nur den
Vergleich zweier Gruppen mit sämtlichen nichtleeren Teilmengen.
Der historische Gruppenrekonstruktionssatz wird dort J.~Shafer,
\emph{Note on power semigroups}, Math.\ Japon.\ 12 (1967), S.~32,
zugeschrieben; diese Originalarbeit wurde für den Leitfaden
nicht unmittelbar ausgewertet.
\url{https://arxiv.org/abs/2606.01917}.

\end{thebibliography}

\noindent
Die Quellen dienen der Begriffs- und Ergebniskontrolle. Mehrere
zusammengehörige Standardaussagen sind hier jeweils zu einem Eintrag
gebündelt; Formulierungen und Nummerierung sind eigenständig.
Die Auswahl enthält weder historische Prioritätsansprüche noch
eine Behauptung vollständiger Erfassung der Gruppenliteratur.
